\documentclass[12pt,a4paper]{letter}
\usepackage[utf8]{inputenc}
\usepackage[T1]{fontenc}
\usepackage{lmodern}
\usepackage[margin=2.5cm]{geometry}

\signature{Gang Zhang\\Independent Researcher\\unicome@gmail.com}
\address{}
\date{March 2026}

\begin{document}
\begin{letter}{The Editors\\Classical and Quantum Gravity}

\opening{Dear Editors,}

Please consider our manuscript, ``Dimensional Selection in Finite
Causal-Set Link Dynamics: Evidence for an Asymmetric $4{\to}5$
Barrier,'' for publication in \emph{Classical and Quantum Gravity}.

This paper studies dimensional competition in finite Lorentzian-like
causal-set ensembles generated from flat-spacetime sprinklings in
dimensions $(1{+}1)$ through $(4{+}1)$.  We assign each family a score
that combines combinatorial entropy with a causal action term, allowing
us to compare the entropic preference for higher-dimensional families
against dynamical penalties arising from causal structure.

The main result is a sharp contrast between two action sectors.  We find
that the link sector of the Benincasa--Dowker action,
$S_{\mathrm{link}} = N - 2C_0$, favors a $(3{+}1)$-dimensional family
within a finite coupling window that remains visible across the tested
sizes and generators.  By contrast, the standard $d{=}4$ BDG action
systematically favors $(4{+}1)$ dimensions over the parameter range
studied.  We further introduce a dimensionless boundary ratio,
$\Xi_{d\to d+1}$, which provides a compact diagnostic of the
action-density separation per unit entropy-density separation between
adjacent dimensions.  Across the tested generators, the $4{\to}5$
boundary is consistently more costly than lower-dimensional transitions,
providing a simple explanation for the appearance of a finite
$(3{+}1)$-dimensional window in the link sector.

We believe this manuscript is well suited to \emph{Classical and Quantum
Gravity} because it addresses causal-set dynamics, discrete spacetime
structure, and the emergence of effective dimension in quantum gravity.
The paper focuses on a finite pre-geometric regime that is complementary
to the better-known continuum-curvature role of the full BDG action, and
it may be of interest to readers working on causal-set theory,
dimensionality in quantum gravity, and discrete approaches to spacetime.

The manuscript is original, has not been published previously, and is not
under consideration elsewhere.  A Zenodo archive containing the numerical
outputs and implementation materials supporting reproducibility is cited
in the manuscript's Data and Code Availability statement.

\closing{Thank you for your consideration.\\[1em]Sincerely,}

\end{letter}
\end{document}
